\documentclass[12pt,a4paper]{article}

\usepackage[utf8]{inputenc}
\usepackage[T2A]{fontenc}
\usepackage{lmodern}
\usepackage[russian]{babel}
\usepackage{amsmath, amssymb, amsthm}
\usepackage{tcolorbox}
\usepackage{enumitem}
\usepackage{array, booktabs}
\usepackage{multicol}
\usepackage[a4paper, margin=1in]{geometry}
\usepackage{fancyhdr}
\usepackage{hyperref}
\usepackage{cleveref}
\usepackage{graphicx}
\usepackage{caption}
\usepackage{footmisc}
\usepackage{xcolor}

\geometry{top=2.5cm, bottom=2.5cm, left=2.5cm, right=2.5cm, headheight=15pt}
\hypersetup{colorlinks=true, linkcolor=blue, filecolor=magenta, urlcolor=cyan, pdfpagemode=FullScreen}

\pagestyle{fancy}
\fancyhf{}
\lhead{\footnotesize \textit{Математические основы теории информации Шеннона и понятие энтропии}}
\rhead{\thepage}

% Custom colors for tcolorbox
\tcbuselibrary{theorems, skins, breakable}
\definecolor{thmcolor}{RGB}{200,230,255}
\definecolor{defcolor}{RGB}{230,250,220}
\definecolor{excolor}{RGB}{255,240,210}

% Theorem environment
\newtheorem{theorem}{Теорема}[section]
\newenvironment{thmbox}[1][]{%
  \begin{tcolorbox}[enhanced, breakable, colback=thmcolor, colframe=blue!50!black, title=#1, fonttitle=\bfseries, arc=3pt, boxrule=0.5pt]
}{\end{tcolorbox}}

% Definition environment
\newtheorem{definition}{Определение}[section]
\newenvironment{defbox}[1][]{%
  \begin{tcolorbox}[enhanced, breakable, colback=defcolor, colframe=green!50!black, title=#1, fonttitle=\bfseries, arc=3pt, boxrule=0.5pt]
}{\end{tcolorbox}}

% Lemma environment
\newtheorem{lemma}{Лемма}[section]
\newenvironment{lembox}[1][]{%
  \begin{tcolorbox}[enhanced, breakable, colback=yellow!10, colframe=orange!50!black, title=#1, fonttitle=\bfseries, arc=3pt, boxrule=0.5pt]
}{\end{tcolorbox}}

% Proposition environment
\newtheorem{proposition}{Утверждение}[section]
\newenvironment{propbox}[1][]{%
  \begin{tcolorbox}[enhanced, breakable, colback=purple!5, colframe=purple!50!black, title=#1, fonttitle=\bfseries, arc=3pt, boxrule=0.5pt]
}{\end{tcolorbox}}

% Corollary environment
\newtheorem{corollary}{Следствие}[section]
\newenvironment{corbox}[1][]{%
  \begin{tcolorbox}[enhanced, breakable, colback=cyan!5, colframe=cyan!50!black, title=#1, fonttitle=\bfseries, arc=3pt, boxrule=0.5pt]
}{\end{tcolorbox}}

% Example environment
\newtheorem{example}{Пример}[section]
\newenvironment{exbox}[1][]{%
  \begin{tcolorbox}[enhanced, breakable, colback=excolor, colframe=brown!50!black, title=#1, fonttitle=\bfseries, arc=3pt, boxrule=0.5pt]
}{\end{tcolorbox}}

% Remark environment
\newtheorem{remark}{Замечание}[section]
\newenvironment{rembox}[1][]{%
  \begin{tcolorbox}[enhanced, breakable, colback=gray!5, colframe=gray!50, title=#1, fonttitle=\bfseries, arc=3pt, boxrule=0.5pt]
}{\end{tcolorbox}}

% Customize enumerate and itemize
\setlist[enumerate]{leftmargin=*, itemsep=2pt, topsep=4pt}
\setlist[itemize]{leftmargin=*, itemsep=2pt, topsep=4pt}

% Custom table environment with caption, placement, and column spec
\newenvironment{customtable}[3][htbp]{%
  % #1 = optional: placement (default: htbp)
  % #2 = mandatory: column specification (e.g., {ccc}, {l|c|r})
  % #3 = mandatory: table caption (goes above the table)
  \begin{table}[#1]
  \centering
  \renewcommand{\arraystretch}{1.2}
  \caption{#3}
  \begin{tabular}{#2}
}{%
  \end{tabular}
  \end{table}
}

\title{Математические основы теории информации Шеннона и понятие энтропии}

% ============================================================
% DOCUMENT STRUCTURE GUIDELINES
% ============================================================
%
% === 1. DOCUMENT STRUCTURE & FLOW ===
% - Use \section{}, \subsection{}, \subsubsection{} to mirror lecture flow.
% - Limit depth: section → subsection → subsubsection (avoid deeper).
% - Start each section with a 1–2 sentence summary of what follows.
% - Example:
%   \section{Introduction}
%   This section introduces the core concepts of vector spaces...

% === 2. CONTENT ACCURACY ===
% - Faithfully represent the lecture: definitions, theorems, examples, logic.
% - Do NOT add external content unless noted (e.g., in footnotes).
% - Paraphrase for clarity; use direct quotes sparingly.

% === 3. MATHEMATICAL TYPESETTING ===
% - Always use math mode: $x^2$, \[ \int f(x) dx \], or \begin{equation}.
% - Number equations only if referenced later (use \label{} and \cref{}).
% - Use semantic commands: \mathbb{R}, \mathcal{L}, \vec{v}, \nabla, etc.
% - Example:
%   \begin{equation}\label{eq:euler}
%   e^{i\pi} + 1 = 0
%   \end{equation}
%   As shown in \cref{eq:euler}, ...

% === 4. CUSTOM ENVIRONMENTS (USE AS INTENDED) ===
% - \begin{defbox}[Title]     → For definitions (green-tinted).
% - \begin{thmbox}[Title]     → For theorems/lemmas/propositions (blue-tinted).
% - \begin{exbox}[Title]      → For examples (beige-tinted).
% - \begin{rembox}[Title]     → For remarks/notes (gray-tinted).
%
% Rules:
% - Always include a title: \begin{thmbox}[Pythagorean Theorem]}
% - Keep content concise; avoid long paragraphs inside boxes.
% - Do NOT nest boxes.
% - For formal numbering, wrap amsthm environments inside:
%   \begin{thmbox}[Mean Value Theorem]
%   \begin{theorem}
%   Let $f$ be continuous on $[a,b]$...
%   \end{theorem}
%   \end{thmbox}

% === 5. LISTS AND ENUMERATIONS ===
% - Use itemize for unordered key points:
%   \begin{itemize}
%     \item Point A
%     \item Point B
%   \end{itemize}
% - Use enumerate for step-by-step procedures:
%   \begin{enumerate}
%     \item Step 1: Define domain.
%     \item Step 2: Compute derivative.
%   \end{enumerate}
% - Avoid nesting >2 levels.

% === 6. TABLES (USE customtable) ===
% - Syntax:
%   \begin{customtable}[placement]{column-spec}{Caption}
%   \toprule
%   Col1 & Col2 & Col3 \\
%   \midrule
%   Data & Data & Data \\
%   \bottomrule
%   \end{customtable}
% - Example:
%   \begin{customtable}[h]{lcc}{Experimental Results}
%   \toprule
%   Trial & Input & Output \\
%   \midrule
%   1 & 10 & 15 \\
%   2 & 20 & 30 \\
%   \bottomrule
%   \end{customtable}
% - Use booktabs rules (\toprule, \midrule, \bottomrule).
% - No vertical lines; align numbers by decimal if needed.


\begin{document}

\maketitle
\begin{abstract}Лекция рассматривает концепции теории информации, введённые Клодом Шенноном, с акцентом на количественную оценку информации и неопределённости. Основной целью является объяснение связи между вероятностью событий и количеством полученной информации. Лектор приводит примеры вероятностных событий для иллюстрации: чем менее вероятно событие, тем больше информации оно несёт. Далее вводится логарифмическая формула, показывающая, что информация определяется логарифмом обратной вероятности. Важное понятие энтропии определяется как среднее значение информации о системе, базируясь на её вероятностных состояниях. Энтропия достигает минимума (0) при полной определённости системы и максимума при равномерном распределении вероятностей по всем возможным состояниям. Таким образом, энтропия служит мерой неопределённости в системе.\end{abstract}

\section{Введение в концепции информации и вероятности}

В данном разделе представлена общая концепция теории информации, основанная на вероятностных событиях, и ключевой вопрос о связи количества передаваемой информации с вероятностью события. Лектор опирается на примеры вероятностных событий, чтобы показать, как они могут нести разное количество информации.

\begin{defbox}[Информация и вероятность]
\textbf{Информация} – это мера уменьшения неопределённости, когда происходит некоторое событие. Связь между вероятностью события $p$ и информацией $I$, полученной при его осуществлении, выражается следующим соотношением:
\[
I = \log \frac{1}{p}.
\]
Чем менее вероятно событие, тем больше информации мы получаем, когда оно происходит.
\end{defbox}

\subsection{Примеры вероятностных событий}

Рассмотрим два вероятностных события, как это было предложено Клодом Шенноном, и как это позже развивал Герман Хакен, работая в области синергетики (теории самоорганизации сложных систем).

\begin{enumerate}
    \item Первое вероятностное событие: профессор $M$ приходит на лекцию.
    \item Второе вероятностное событие: профессор $M$ ругает студента.
\end{enumerate}

Теперь зададим вопрос: какое из этих событий приносит больше информации?

\begin{exbox}[Сравнение двух событий]
Пусть событие «профессор $M$ приходит на лекцию» имеет более высокую вероятность (то есть это привычное и ожидаемое событие). В таком случае, согласно формуле информации $I = \log \frac{1}{p}$, количество информации, полученное при наступлении этого события, будет относительно небольшим, так как $p$ велико.

С другой стороны, событие «профессор $M$ ругает студента» может быть редким и маловероятным, а значит, его вероятность $p$ мала. В этом случае информация $I$, полученная при его наступлении, будет значительно больше.
\end{exbox}

\begin{rembox}[Ключевое замечание]\label{remark:event_information}
Следовательно, чем более неожиданное событие (тем ниже его вероятность), тем больше информации оно несёт. Именно это соотношение лежит в основе количественной оценки информации в теории Шеннона.
\end{rembox}

\subsection{Логика оценки информации}

Более формально, информация — это мера неожиданности события. Если событие происходит часто, то оно лишь незначительно уменьшает неопределённость относительно системы, и, следовательно, информации в таком событии мало.

Напротив, если событие крайне редкое, то оно значительно уменьшает неопределённость, и его информация будет высокой. Именно поэтому редкие события несут больше информации, чем частые события. Таким образом, информация и вероятность события связаны обратно пропорционально.

\begin{rembox}[Контекст синергетики]
Традиционно эта концепция информации рассматривается в контексте синергетики, поскольку синергетика изучает самоорганизацию сложных систем и требует понимания того, как информация возникает из маловероятных событий, которые способны изменить поведение системы.
\end{rembox}

Любая система, будь то социальная, физическая или биологическая, содержит различные вероятностные состояния, и понимание того, как события в этих системах связаны с информацией, является ключевым для анализа их поведения и эволюции.
\subsubsection{Соотношение между вероятностью и интересом к событию}

Рассмотрим простой гипотетический пример: у нас есть два события. Первое событие — профессор задаёт вопрос студенту в аудитории. Второе событие — профессор задаёт вопрос студенту на улице. Какой из этих сценариев может быть более интересен для масс-медиа?

\begin{itemize}
    \item Первый сценарий достаточно обычен: профессора часто взаимодействуют со студентами в аудитории.
    \item Второй сценарий, напротив, необычен — ведь редко можно увидеть, чтобы профессор внезапно начал опрашивать студентов на улице, вне формальной учебной обстановки.
\end{itemize}

\begin{rembox}[Интересность события]
Ожидается, что событие, обладающее меньшей вероятностью, обычно воспринимается как более интересное и несёт в себе больше информации. Это следует из того, что маловероятные события содержат в себе большую степень неожиданности.
\end{rembox}

Таким образом, второй сценарий привлекает больше внимания и несёт в себе больше информации, поскольку вероятность такого события значительно ниже.

\subsubsection{Суммарная информация для независимых событий}

Перейдём к следующему логическому шагу. Информация, как мы уже убедились, является функцией от обратной вероятности. Обозначим информацию от события $A$ как $I(A)$.

\begin{defbox}[Функция информации]
Информация является функцией от обратной вероятности события. Формально, для события $A$ с вероятностью $P(A)$, информация $I(A)$ может быть представлена в виде:
\[
I(A) = f\left( \frac{1}{P(A)} \right),
\]
где $f$ — некоторая монотонная функция.
\end{defbox}

Теперь предположим, что у нас есть два независимых* случайных события $A$ и $B$. В силу независимости их совместная вероятность равна произведению вероятностей отдельных событий:
\[
P(A \cap B) = P(A) \cdot P(B).
\]

\begin{rembox}[Независимость событий]
Независимость событий означает, что наступление одного события не влияет на наступление другого. В статистике это выражается тем, что $P(A \cap B) = P(A) \cdot P(B)$.
\end{rembox}

Возникает вопрос: какова суммарная информация об обоих событиях $A$ и $B$, если они независимы? Интуиция говорит, что информация суммируется, ведь каждое событие вносит свой вклад в общую неопределённость системы (или уменьшает её).

\begin{thmbox}[Сумма информации независимых событий]
Пусть $A$ и $B$ — два независимых события. Суммарная информация о двух таких событиях равна сумме информации об отдельном первом и отдельном втором событии:
\[
I(A \cap B) = I(A) + I(B).
\]
\end{thmbox}

Это второй логический шаг Шеннона: информация для независимых событий складывается. Доказательство этого утверждения исходит из логарифмических свойств вероятностей.

\begin{rembox}[Интуиция логарифмической функции]
Логарифмы являются естественным выбором для измерения информации, потому что они переводят произведение в сумму. Именно на этом принципе базируется определение информации для независимых событий.
\end{rembox}

В следующих разделах мы более подробно рассмотрим, как логарифмические функции используются для определения информации и почему именно они были выбраны Шенноном как мера информации. 

Сейчас важно понять, что общая информация о системе, состоящей из множества независимых событий, равна сумме информации об отдельных событиях. Эта простая, но важная идея лежит в основе всей теории информации.
\section{Свойства информации и логарифмическая функция}

В данной части лекции обсуждается важное свойство информации: для независимых событий суммарная информация равна сумме информации отдельных событий. Этот факт приводит к необходимости использования логарифмической функции как единственной, которая удовлетворяет данному свойству.

\begin{thmbox}[Свойства информации]
\begin{theorem}
Если два события $A$ и $B$ независимы, то информация о наступлении обоих событий $I(A \cap B)$ равна сумме информации о каждом из них: 
\[
I(A \cap B) = I(A) + I(B).
\]
\end{theorem}
\end{thmbox}

Единственная функция, которая удовлетворяет этому свойству, — логарифмическая функция. Это утверждение интуитивно понятно, но, тем не менее, существует строгое математическое доказательство того, что именно логарифм соответствует этим условиям. 

\subsection{Логарифмическое выражение информации}

Важно подчеркнуть, что информация о событии связана с обратной вероятностью этого события. Формально:
\[
I = \log\left(\frac{1}{p}\right),
\]
где $p$ — вероятность события.

При этом традиционно используется натуральный логарифм (по основанию $e$), однако при использовании логарифма по основанию 2 получаем более привычные величины информации, измеряемые в битах.

\begin{rembox}[Основание логарифма]
Выбор основания логарифма влияет на единицы измерения информации:
\begin{itemize}
    \item Если основание — $e$, информация измеряется в натах.
    \item Если основание — 2, получаем информацию в битах.
\end{itemize}
\end{rembox}

\section{Определение энтропии}

Понятие энтропии тесно связано с вероятностными состояниями системы. Рассмотрим систему, которая может находиться в одном из $n$ состояний. Вероятность нахождения системы в $i$-м состоянии обозначим как $p_i$.

\begin{defbox}[Энтропия системы]
Энтропия системы — это средняя информация, связанная с неопределённостью её состояния. Выражается она следующим образом:
\[
H = -\sum_{i=1}^{n} p_i \log(p_i).
\]
\end{defbox}

Здесь $H$ обозначает энтропию (обычно в битах, если основание логарифма равно 2), а $p_i$ — вероятность нахождения системы в данном состоянии $i$.

\begin{rembox}[Физический смысл энтропии]
Энтропия измеряет степень неопределённости (или случайности) системы: 
\begin{itemize}
    \item Если система всегда находится в одном и том же состоянии (т.е. одно состояние имеет вероятность 1, а остальные — 0), энтропия равна $0$.
    \item Если вероятности всех состояний равномерно распределены (например, $p_i = \frac{1}{n}$ для всех $i$), энтропия достигает максимального значения.
\end{itemize}
\end{rembox}
\subsection{Интерпретация энтропии}

Энтропия, введённая Клодом Шенноном, является чрезвычайно интересной мерой оценки информации. По сути, энтропия представляет собой среднее количество информации о системе, которую мы получаем при наблюдении её состояния.

\begin{rembox}[Две интерпретации энтропии]
Существует два способа понимания или интерпретации концепции энтропии:
\begin{enumerate}
    \item \textbf{Средняя неопределённость состояния:} 
    \begin{itemize}
        \item Энтропия описывает, насколько неопределённо, в среднем, текущее состояние системы до того, как мы наблюдаем этот состояние.
        \item Чем выше энтропия, тем больше неопределённости и тем труднее предсказать состояние системы.
    \end{itemize}
    \item \textbf{Средний объём информации:}
    \begin{itemize}
        \item С другой стороны, энтропию можно рассматривать как среднее количество информации, которое даёт нам наблюдение состояния системы.
        \item Если состояние является маловероятным, то полученная информация велика, а если оно очень вероятно, то информация, получаемая при наблюдении, мала.
    \end{itemize}
\end{enumerate}
\end{rembox}

Таким образом, энтропия отражает как уровень неопределённости, так и среднюю информацию, извлекаемую из системы.

\subsection{Алгебраические свойства энтропии}

Далее, рассмотрим несколько важных свойств энтропии, которые помогают глубже понять её поведение при различных распределениях вероятностей.

\begin{thmbox}[Свойства энтропии]
\begin{theorem}
Пусть $p = (p_1, p_2, \dots, p_n)$ — распределение вероятностей для конечной дискретной системы. Тогда энтропия $H(p)$ обладает следующими свойствами:
\begin{enumerate}
    \item \textbf{Неотрицательность}:
    \[
    H(p) \geq 0.
    \]
    Энтропия всегда неотрицательна, так как каждое слагаемое суммы $\left(-p_i \log p_i\right)$ неотрицательно для всех $p_i \in [0,1]$.
    
    \item \textbf{Минимальное значение энтропии}:
    \[
    H(p) = 0,
    \]
    если и только если всё распределение сосредоточено на одном состоянии (например, $p_1=1, p_2=0, \dots, p_n=0$). В этом случае нет неопределённости, и наблюдение не даёт новой информации.
    
    \item \textbf{Максимальное значение энтропии}:
    \[
    H(p) = \log(n)
    \]
    когда все вероятности равны: $p_i = \frac{1}{n}$ для всех $i$. В этом случае неопределённость максимальна.
    
    \item \textbf{Аддитивность для независимых систем}:
    Если две системы $A$ и $B$ независимы, с распределениями $p(A)$ и $p(B)$, то энтропия их совместной системы равна:
    \[
    H(A, B) = H(A) + H(B).
    \]
    \item \textbf{Вогнутость}:
    Функция энтропии является вогнутой функцией распределения вероятностей, то есть:
    \[
    H(\lambda p_1 + (1-\lambda) p_2) \geq \lambda H(p_1) + (1 - \lambda) H(p_2)
    \]
    для любых распределений $p_1, p_2$ и $0 \leq \lambda \leq 1$.
\end{enumerate}
\end{theorem}
\end{thmbox}

Эти свойства помогают математически работать с энтропией и предсказывать, как она изменяется в зависимости от изменения распределения вероятностей.

\begin{rembox}[Практическое значение вогнутости энтропии]
Вогнутость энтропии означает, что смешивание двух различных распределений вероятностей даёт энтропию, как минимум не меньшую, чем средневзвешенную энтропию этих распределений. Это важно для оценки информации в системах с неоднородными источниками сигналов.
\end{rembox}
\subsection{Канальная энтропия и условная энтропия}

В теории информации Шеннона важными понятиями являются условная энтропия и взаимная информация. Условная энтропия позволяет оценить неопределённость одной случайной величины при условии, что известна другая величина.

\begin{defbox}[Условная энтропия]
\begin{definition}
Условная энтропия случайной величины $X$ относительно другой случайной величины $Y$ определяется как:
\[
H(X|Y) = - \sum_{y \in \mathcal{Y}} p(y) \sum_{x \in \mathcal{X}} p(x|y) \log p(x|y),
\]
где $p(x|y)$ — условная вероятность $X$ при фиксированном $Y=y$, а $p(y)$ — маргинальная вероятность $Y=y$.
\end{definition}
\end{defbox}

\begin{rembox}[Смысл условной энтропии]
Условная энтропия $H(X|Y)$ измеряет среднюю неопределённость случайной величины $X$, учитывая, что информация о $Y$ уже известна. 
\end{rembox}

Таким образом, условная энтропия описывает, сколько информации нам ещё нужно для определения $X$, если мы уже что-то знаем о $Y$.

\subsection{Связь между полной энтропией и условной энтропией}

Существует тесная связь между совместной энтропией двух переменных, их индивидуальными энтропиями и условной энтропией. Это позволяет лучше понять структуру систем и источников информации.

\begin{thmbox}[Основные равенства для совместной и условной энтропии]
\begin{theorem}
Пусть $X$ и $Y$ — две дискретные случайные величины с распределением вероятностей $p(x,y)$. Тогда справедливы следующие равенства:
\begin{enumerate}
    \item \textbf{Совместная энтропия}:
    \[
    H(X,Y) = H(X) + H(Y|X) = H(Y) + H(X|Y),
    \]
    где $H(X,Y)$ — совместная энтропия обеих величин.
    
    \item \textbf{Неравенство условной энтропии}:
    \[
    H(X|Y) \leq H(X).
    \]
    Условная энтропия всегда не больше, чем полная энтропия $X$, так как информация о $Y$ уменьшает неопределённость относительно $X$.
\end{enumerate}
\end{theorem}
\end{thmbox}

\begin{rembox}[Интуитивное объяснение]
Совместная энтропия $H(X,Y)$ — это мера неопределённости обеих величин вместе. Условная энтропия показывает, как неопределённость одной величины уменьшается, если мы знаем другую. Если $X$ и $Y$ независимы, условная энтропия $H(X|Y)$ равна полной энтропии $H(X)$, ведь $Y$ не даёт никакой дополнительной информации о $X$.
\end{rembox}

\subsection{Взаимная информация}

Теперь мы можем определить важную величину, называемую \textit{взаимной информацией}, которая измеряет объём информации, который одна случайная величина даёт о другой.

\begin{defbox}[Взаимная информация]
\begin{definition}
Взаимная информация между случайными величинами $X$ и $Y$ определяется следующим образом:
\[
I(X; Y) = H(X) - H(X|Y) = H(Y) - H(Y|X).
\]
\end{definition}
\end{defbox}

\begin{rembox}[Смысл взаимной информации]
Взаимная информация $I(X; Y)$ показывает, насколько уменьшится неопределённость относительно $X$, если мы узнаем $Y$, и наоборот. Чем больше $I(X; Y)$, тем больше зависимость между $X$ и $Y$.
\end{rembox}

Взаимная информация — это мера зависимости между двумя переменными: она равна нулю, если переменные независимы, и увеличивается при наличии связи между ними.

\subsection{Формулы для взаимной информации}

Взаимная информация имеет несколько эквивалентных представлений, которые полезны при расчётах.

\begin{thmbox}[Формулы взаимной информации]
\begin{theorem}
Взаимная информация $I(X; Y)$ может быть выражена следующими эквивалентными способами:
\begin{enumerate}
    \item Через совместную энтропию:
    \[
    I(X; Y) = H(X) + H(Y) - H(X, Y).
    \]
    \item Через условные вероятности:
    \[
    I(X; Y) = \sum_{x \in \mathcal{X}} \sum_{y \in \mathcal{Y}} p(x,y) \log\left(\frac{p(x,y)}{p(x)p(y)}\right),
    \]
    где $p(x,y)$ — совместное распределение $X$ и $Y$, а $p(x)$ и $p(y)$ — маргинальные распределения.
\end{enumerate}
\end{theorem}
\end{thmbox}

\begin{rembox}[Объяснение в терминах вероятности]
Если $X$ и $Y$ независимы, то $p(x,y) = p(x)p(y)$, и взаимная информация $I(X; Y)$ равна нулю. Если существует сильная зависимость между $X$ и $Y$, то $p(x,y)$ значительно отличается от произведения $p(x)p(y)$, что увеличивает $I(X; Y)$.
\end{rembox}

Эти формулы показывают, что взаимная информация тесно связана и с энтропией отдельных величин, и с их совместной энтропией.

\begin{exbox}[Пример расчёта взаимной информации]
Рассмотрим случайные величины $X$ и $Y$, каждая из которых может принимать значения $0$ и $1$ с равными вероятностями $p(X=0) = 0.5$, $p(X=1) = 0.5$, и аналогично для $Y$. Пусть они полностью независимы. Тогда
\[
p(x,y) = p(x) \cdot p(y).
\]
В таком случае:
\[
I(X; Y) = \sum_{x,y} p(x,y) \log \left( \frac{p(x,y)}{p(x)p(y)} \right) = 0.
\]

Если же $X$ и $Y$ связаны, например $Y = X$, то $p(x,y) = p(x)$, и тогда:
\[
I(X; Y) = H(X),
\]
так как знание $Y$ полностью определяет $X$.
\end{exbox}

\subsection{Аддитивное свойство взаимной информации}

Важным аспектом взаимной информации является её аддитивность при независимых системах.

\begin{thmbox}[Аддитивность взаимной информации]
\begin{theorem}
Пусть у нас есть две пары независимых случайных величин $(X_1, Y_1)$ и $(X_2, Y_2)$. Тогда взаимная информация для объединённой системы будет равна сумме взаимных информаций каждой отдельной пары:
\[
I(X_1, X_2; Y_1, Y_2) = I(X_1; Y_1) + I(X_2; Y_2).
\]
\end{theorem}
\end{thmbox}

Это свойство полезно для анализа сложных систем, которые можно разбить на отдельные подсистемы, обрабатывая их взаимную информацию по частям.

\begin{rembox}[Практическое применение аддитивности]
Аддитивность позволяет разбивать большие системы на подмножества и анализировать их информационные зависимости по частям, что упрощает расчёты и понимание.
\end{rembox}

Далее мы рассмотрим важный результат Шеннона — теорему о передаче информации через канал с шумом, которая использует понятия энтропии и взаимной информации для описания предельных возможностей каналов связи.
\subsection{Неопределённость и количество информации}

Рассмотрим систему, которая может находиться в нескольких допустимых состояниях, каждое из которых имеет определённую вероятность $p_i$. После наблюдения системы и фиксации её в некотором конкретном состоянии — например, в первом, — становится ясно, что мы получили определённую информацию об этой системе. Вопрос заключается в том, сколько информации мы получили.

\begin{defbox}[Количество информации в наблюдаемом событии]
Количество информации, полученное при наблюдении системы в конкретном состоянии, определяется как:
\[
I_i = -\log_2 p_i,
\]
где $p_i$ — вероятность того, что система находится в этом состоянии.
\end{defbox}

\begin{rembox}[Интерпретация количества информации]
Чем менее вероятно событие (меньше $p_i$), тем больше информации оно несёт при наступлении. Например, если событие крайне маловероятно, его наступление передаёт очень большой объём информации о системе.
\end{rembox}

Таким образом, количество информации, которое мы получаем, зависит от того, насколько неожиданным было наблюдаемое состояние.

С другой стороны, если мы знаем только вероятностное распределение состояний системы, но не наблюдаем её текущее состояние, наша информация о системе остаётся неполной. Наша неопределённость относительно состояния системы в таком случае характеризуется некоторой мерой. Возникает естественный вопрос: как измерить эту неопределённость?

\subsection{Энтропия как мера неопределённости}

Энтропия системы — это мера неопределённости относительно её текущего состояния при известном распределении вероятностей. Если бы мы знали точное состояние системы, энтропия была бы равна нулю, поскольку неопределённость отсутствует. Если же система полностью неизвестна и все состояния равновероятны, энтропия достигает максимума.

\begin{defbox}[Энтропия системы]
Энтропия $H$ системы с вероятностным распределением $p_1, p_2, \dots, p_n$ определяется как:
\[
H = - \sum_{i=1}^n p_i \log_2 p_i.
\]
\end{defbox}

Энтропия служит количественной мерой неопределённости системы. Если все состояния равновероятны (например, при $n$ состояниях $p_i = \frac{1}{n}$), то энтропия достигает своего максимума:
\[
H_{\text{max}} = \log_2 n.
\]

\begin{exbox}[Пример]
Допустим, система имеет три состояния с вероятностями $p_1 = 0.5$, $p_2 = 0.3$, $p_3 = 0.2$. Тогда энтропия системы будет равна:
\[
H = - \left( 0.5 \log_2 0.5 + 0.3 \log_2 0.3 + 0.2 \log_2 0.2 \right).
\]
\end{exbox}

\begin{rembox}[Интуитивное понимание энтропии]
Энтропия максимальна, когда система максимально неопределённа, то есть когда все состояния равновероятны. Если одно из состояний доминирует (имеет вероятность, близкую к 1), энтропия стремится к нулю, так как неопределённость уменьшается.
\end{rembox}

Энтропия — это фундаментальное понятие теории информации. Она не только измеряет неопределённость, но и формирует основу для понимания пределов передачи информации в системах связи и хранения данных. В следующем разделе мы рассмотрим, как концепция энтропии используется в теории каналов связи, сформулированной Шенноном.
\subsection{Минимум энтропии и его достижение}

Мы можем провести строгий анализ выражения для энтропии и установить, что её минимальное значение равно нулю. Важно подчеркнуть, что энтропия достигает своего минимума в следующем случае: когда вероятность одного из состояний равна 1, а все остальные вероятности равны 0. Рассмотрим этот случай подробнее.

\begin{exbox}[Пример минимальной энтропии]
Предположим, что вероятность первого состояния $p_1 = 1$, а остальные $p_2, p_3, \dots, p_n = 0$. Тогда энтропия системы:
\[
H = - \left( 1 \cdot \log_2 1 + 0 \cdot \log_2 0 + \dots \right).
\]
Поскольку $\log_2 1 = 0$, результат равен:
\[
H = 0.
\]
\end{exbox}

Таким образом, если мы знаем, что система находится точно в одном состоянии, возможность неопределённости исчезает, и энтропия системы становится равной нулю.

\begin{rembox}[Нулевая энтропия]
Нулевая энтропия означает, что мы знаем абсолютно всё о системе. В этом случае мера неопределённости равна нулю, поскольку никаких неизвестных компонентов в системе не остаётся.
\end{rembox}

С другой стороны, если мы совершенно ничего не знаем о состоянии системы, или если все состояния равновероятны, энтропия достигает своего максимального значения, как мы уже обсудили ранее. Это отражает максимальную неопределённость системы.

\subsection{Заключение}

Подводя итог, энтропия является ключевой величиной в теории информации. Она служит мерой неопределённости относительно состояния системы и помогает количественно оценить объём информации, который может быть передан через канал связи. Когда одна вероятность состояния резко доминирует, энтропия падает до нуля, что соответствует полному знанию о системе. В противоположном случае, когда все события равновероятны, энтропия максимальна, отражая максимальную неопределённость.

Концепция энтропии, как мы видим, лежит в основе понимания пределов передачи информации и является фундаментальным инструментом для анализа систем хранения данных и передачи сигналов. Дальнейшие исследования и разработки в области теории информации строятся на этом фундаменте, заложенном Клодом Шенноном.

\end{document}
